\chapter*{Conclusion générale et perspectives}
\addcontentsline{toc}{chapter}{Conclusion générale et perspectives}

Le stage a permis de concevoir et d'implémenter \textbf{SmartTrans-AI}, une application de gestion de transport centrée sur la mission, les rôles et un assistant ancré sur les données. Du point de vue pédagogique, le travail a mobilisé l'état de l'art, la spécification UML, la conception d'une API REST sécurisée, une interface selon le profil utilisateur, une première intégration d'IA contrainte par le SI, ainsi que la conduite d'un projet itératif.

Du point de vue de l'organisme, le livrable offre une base opérationnelle pour le département informatique de S. WATRANS : planification, suivi, alertes d'expiration et traçabilité. Les perspectives immédiates concernent le renforcement des tests métier, le déploiement sécurisé, et d'éventuels modules (notifications, pièces jointes, interfaçage temps réel) si le besoin s'en confirme.

Ce rapport, tenu dans une fourchette d'une quarantaine de pages (limite ENSIASD : soixante-dix pages), constitue la trace écrite de cette réalisation.
